\section{Generalisable Feed-Forward 3D Reconstruction}
\label{sec:generalisable}

\subsection{Pose-Free and Foundation Models}

A new generation of methods aims for \emph{pose-free} generalizable reconstruction.
DUSt3R~\citep{wang2024dust3r} predicts dense pointmaps from arbitrary image pairs
without camera calibration, then aligns them via global optimisation.
MASt3R~\citep{leroy2024mast3r} adds matching capabilities to DUSt3R, enabling
accurate localisation.
MonST3R~\citep{zhang2024monst3r} extends DUSt3R to dynamic video.
VGGT~\citep{wang2025vggt} is a visual geometry grounded deep structured tokens
model for feed-forward 3-D structure prediction.

\subsection{Large Reconstruction Models (LRM)}
\label{sec:lrm}

The Large Reconstruction Model (LRM)~\citep{hong2023lrm} paradigm represents
a qualitative shift from per-scene optimisation to \emph{feed-forward inference}:
a single forward pass of a large transformer, pretrained on hundreds of thousands
of 3-D assets, produces a NeRF or Gaussian scene from a single image in under
5 seconds.
This is possible because the model amortises the geometric reasoning of NeRF
optimisation into its weights---a striking demonstration that scale can substitute
for per-instance compute.

LRM~\citep{hong2023lrm} itself uses an image encoder (DINO-ViT) followed by
a cross-attention transformer that maps image tokens to triplane tokens, which
are then decoded by a NeRF MLP.
Instant3D~\citep{li2023instant3d} extends this by first generating sparse multi-view
images with a diffusion model, then feeding the result to an LRM-style reconstructor,
decoupling appearance generation from geometry prediction.
OpenLRM~\citep{openlrm2023} reproduces LRM in an open-source setting.

The Gaussian-based branch has proven especially fast.
LGM~\citep{tang2024lgm} generates multi-view images and lifts them to a 3DGS
scene via an asymmetric U-Net, achieving 3-D generation in seconds at
commercial-grade quality.
GRM~\citep{xu2024grm} predicts 3D Gaussians directly from image features using
a sparse-view transformer, operating entirely at the pixel level without any
intermediate NeRF.
Gamba~\citep{he2024gamba} replaces the transformer with a state-space model
(Mamba) for $O(N)$ sequence length scaling, enabling long-context 3-D inference.
GS-LRM~\citep{zhang2024gslrm} combines a transformer backbone with a 3DGS
output head, achieving state-of-the-art single-image reconstruction quality.

Two non-parametric reconstruction models have recently emerged as strong baselines.
DUSt3R~\citep{wang2024dust3r} formulates 3-D reconstruction as point-map
regression: given an image pair, a ViT predicts dense 3-D coordinates for
every pixel in a shared Euclidean frame, requiring no camera poses at test time.
MASt3R~\citep{leroy2024mast3r} extends DUSt3R to feature-level matching,
enabling accurate metric reconstruction from internet-sourced photo pairs.
These models underpin a new generation of unposed feed-forward reconstructors
such as Splatt3R~\citep{smart2024splatt3r}, which produces Gaussian splatting
scenes directly from uncalibrated image pairs.

\paragraph{Implications for World Modeling.}
LRM-style models demonstrate that a single transformer, with sufficient data
and compute, can learn a universal spatial prior---the same capacity required
by world models for instant scene initialisation from single-frame observations.

\subsection{Multi-View Stereo in the Neural Era}

Classical MVS has been reinvigorated by deep learning.
MVSNet~\citep{yao2018mvsnet} introduced cost-volume-based deep MVS.
IterMVS~\citep{wang2022itermvs} and Uni-MVSNet~\citep{peng2022unimvsnet}
improved accuracy and generalisation.
These methods provide the geometry backbone for downstream NVS, and their
cost-volume feature matching is directly incorporated into generalizable NeRF
(MVSNeRF~\citep{chen2021mvsnerf}) and feed-forward Gaussian methods
(MVSplat~\citep{chen2024mvsplat}).

% ─────────────────────────────────────────────────────────────────────────────
